\section{酸碱理论}\label{ux9178ux78b1ux7406ux8bba}

\begin{center}\rule{0.5\linewidth}{0.5pt}\end{center}

\subsection{学习目标}\label{ux5b66ux4e60ux76eeux6807}

\begin{itemize}
\tightlist
\item
  目标1：能用 Brønsted 质子理论识别酸、碱和两性物质，写出共轭酸碱对
\item
  目标2：能正确书写弱酸弱碱的电离平衡表达式，理解 Ka · Kb = Kw 的推导与应用
\item
  目标3：能用简化公式和精确法(假设-验证法)计算弱酸弱碱溶液的 pH
\item
  目标4：能计算两性物质溶液的 pH(NaHCO\textsubscript{3}、NaH\textsubscript{2}PO\textsubscript{4} 等)
\item
  目标5：能设计和计算缓冲溶液，理解 Henderson-Hasselbalch 方程的适用条件
\item
  目标6：了解 Lewis 酸碱理论和 HSAB 软硬分类的基本规则
\item
  目标7：了解溶剂理论中的拉平效应和区分效应
\end{itemize}

\begin{center}\rule{0.5\linewidth}{0.5pt}\end{center}

\subsection{一、为什么需要酸碱理论}\label{ux4e00ux4e3aux4ec0ux4e48ux9700ux8981ux9178ux78b1ux7406ux8bba}

\textbf{酸碱理论的演进}：

{\def\LTcaptype{none} % do not increment counter
\begin{longtable}[]{@{}lclll@{}}
\toprule\noalign{}
理论 & 年代 & 酸的定义 & 碱的定义 & 适用范围 \\
\midrule\noalign{}
\endhead
\bottomrule\noalign{}
\endlastfoot
\textbf{Arrhenius} & 1887 & 电离出 H\textsuperscript{+} & 电离出 OH\textsuperscript{-} & 水溶液 \\
\textbf{Brønsted} & 1923 & 质子供体 & 质子受体 & 质子转移反应 \\
\textbf{溶剂理论} & 1923 & 溶剂阳离子的来源 & 溶剂阴离子的来源 & 非水溶剂 \\
\textbf{Lewis} & 1923 & 电子对受体 & 电子对供体 & 所有酸碱反应 \\
\textbf{HSAB} & 1963 & 软/硬酸分类 & 软/硬碱分类 & 配位化学、有机金属 \\
\end{longtable}
}

\begin{quote}
\textbf{核心认知}：酸碱不是物质的固有属性，而是\textbf{反应角色}。同一物质在不同反应中可以扮演酸或碱---``有酸才有碱，有碱才有碱，酸中有碱，碱可变酸''。
\end{quote}

\begin{center}\rule{0.5\linewidth}{0.5pt}\end{center}

\subsection{二、Brønsted 质子理论}\label{ux4e8cbruxf8nsted-ux8d28ux5b50ux7406ux8bba}

\subsubsection{2.1 基本定义}\label{ux57faux672cux5b9aux4e49}

\textbf{酸}：质子供体(能给出 H\textsuperscript{+} 的物质)
\textbf{碱}：质子受体(能接受 H\textsuperscript{+} 的物质)

\[
\text{酸} \rightleftharpoons \text{碱} + \mathrm{H}^+
\]

\textbf{共轭酸碱对}：酸失去一个质子后变成其共轭碱，碱接受一个质子后变成其共轭酸。

\begin{quote}
\textbf{类比理解}：共轭酸碱就像硬币的两面---``共轭酸碱对是同一物质的两种身份''。
\end{quote}

{\def\LTcaptype{none} % do not increment counter
\begin{longtable}[]{@{}lll@{}}
\toprule\noalign{}
酸 & 共轭碱 & 说明 \\
\midrule\noalign{}
\endhead
\bottomrule\noalign{}
\endlastfoot
HCl & Cl\textsuperscript{-} & 强酸的共轭碱极弱 \\
HAc & Ac\textsuperscript{-} & 弱酸的共轭碱是弱碱 \\
H\textsubscript{2}O & OH\textsuperscript{-} & 水既是酸又是碱 \\
NH\textsubscript{4}\textsuperscript{+} & NH\textsubscript{3} & 铵离子是酸 \\
H\textsubscript{3}O\textsuperscript{+} & H\textsubscript{2}O & 水合氢离子是酸 \\
\end{longtable}
}

\subsubsection{2.2 两性物质}\label{ux4e24ux6027ux7269ux8d28}

既能给出质子又能接受质子的物质：

{\def\LTcaptype{none} % do not increment counter
\begin{longtable}[]{@{}
  >{\raggedright\arraybackslash}p{(\linewidth - 4\tabcolsep) * \real{0.3333}}
  >{\raggedright\arraybackslash}p{(\linewidth - 4\tabcolsep) * \real{0.3333}}
  >{\raggedright\arraybackslash}p{(\linewidth - 4\tabcolsep) * \real{0.3333}}@{}}
\toprule\noalign{}
\begin{minipage}[b]{\linewidth}\raggedright
物质
\end{minipage} & \begin{minipage}[b]{\linewidth}\raggedright
作为酸
\end{minipage} & \begin{minipage}[b]{\linewidth}\raggedright
作为碱
\end{minipage} \\
\midrule\noalign{}
\endhead
\bottomrule\noalign{}
\endlastfoot
\textbf{H\textsubscript{2}O} & H\textsubscript{2}O \箭 OH\textsuperscript{-} + H\textsuperscript{+} & H\textsubscript{2}O + H\textsuperscript{+} \箭 H\textsubscript{3}O\textsuperscript{+} \\
\textbf{HCO\textsubscript{3}\textsuperscript{-}} & HCO\textsubscript{3}\textsuperscript{-} \箭 CO\textsubscript{3}\textsuperscript{2-} + H\textsuperscript{+} & HCO\textsubscript{3}\textsuperscript{-} + H\textsuperscript{+} \箭 H\textsubscript{2}CO\textsubscript{3} \\
\textbf{HPO\textsubscript{4}\textsuperscript{2-}} & HPO\textsubscript{4}\textsuperscript{2-} \箭 PO\textsubscript{4}\textsuperscript{3-} + H\textsuperscript{+} & HPO\textsubscript{4}\textsuperscript{2-} + H\textsuperscript{+} \箭 H\textsubscript{2}PO\textsubscript{4}\textsuperscript{-} \\
\textbf{H\textsubscript{2}O} & NH\textsubscript{3} + H\textsubscript{2}O \箭 NH\textsubscript{4}\textsuperscript{+} + OH\textsuperscript{-}(水是酸) & NH\textsubscript{3} + H\textsubscript{2}O \左箭 NH\textsubscript{4}\textsuperscript{+} + OH\textsuperscript{-}(水是碱) \\
\end{longtable}
}

\begin{quote}
水是最典型的两性物质---``这是 Brønsted 理论最优雅的推论之一''。在不同反应中，水可以扮演酸的角色(如与 NH\textsubscript{3} 反应)，也可以扮演碱的角色(如与 HCl 反应)。
\end{quote}

\subsubsection{2.3 Brønsted 反应的三种类型}\label{bruxf8nsted-ux53cdux5e94ux7684ux4e09ux79cdux7c7bux578b}

{\def\LTcaptype{none} % do not increment counter
\begin{longtable}[]{@{}
  >{\raggedright\arraybackslash}p{(\linewidth - 4\tabcolsep) * \real{0.3333}}
  >{\raggedright\arraybackslash}p{(\linewidth - 4\tabcolsep) * \real{0.3333}}
  >{\raggedright\arraybackslash}p{(\linewidth - 4\tabcolsep) * \real{0.3333}}@{}}
\toprule\noalign{}
\begin{minipage}[b]{\linewidth}\raggedright
类型
\end{minipage} & \begin{minipage}[b]{\linewidth}\raggedright
示例
\end{minipage} & \begin{minipage}[b]{\linewidth}\raggedright
特点
\end{minipage} \\
\midrule\noalign{}
\endhead
\bottomrule\noalign{}
\endlastfoot
\textbf{电离} & HAc + H\textsubscript{2}O \衡 H\textsubscript{3}O\textsuperscript{+} + Ac\textsuperscript{-} & 酸把质子给水 \\
\textbf{自偶电离} & H\textsubscript{2}O + H\textsubscript{2}O \衡 H\textsubscript{3}O\textsuperscript{+} + OH\textsuperscript{-} & 水既给质子又接受质子 \\
\textbf{中和} & H\textsubscript{3}O\textsuperscript{+} + OH\textsuperscript{-} \衡 2H\textsubscript{2}O & 酸和碱直接反应 \\
\end{longtable}
}

\subsubsection{2.4 拉平效应与区分效应}\label{ux62c9ux5e73ux6548ux5e94ux4e0eux533aux5206ux6548ux5e94}

\begin{quote}
\textbf{溶剂理论的独特贡献}：不同溶剂对酸碱强度有不同的''拉平''或''区分''作用。
\end{quote}

\textbf{拉平效应}(Leveling Effect)：在水溶液中，所有比 H\textsubscript{3}O\textsuperscript{+} 强的酸都被''拉平''到 H\textsubscript{3}O\textsuperscript{+} 的水平。

\[
\mathrm{HClO_4 + H_2O \to H_3O^+ + ClO_4^-} \quad (\text{完全电离})
\]
\[
\mathrm{HCl + H_2O \to H_3O^+ + Cl^-} \quad (\text{完全电离})
\]

在水中，HClO\textsubscript{4}、HCl、HNO\textsubscript{3} 都是''强酸''---因为它们都把质子完全转移给了水，生成的都是 H\textsubscript{3}O\textsuperscript{+}。水''拉平''了它们的强度差异。

\textbf{区分效应}(Differentiating Effect)：换一种溶剂(如冰醋酸)，可以区分这些酸的强弱。

\[
\mathrm{HClO_4 + CH_3COOH \rightleftharpoons CH_3COOH_2^+ + ClO_4^-} \quad (\text{完全电离})
\]
\[
\mathrm{HCl + CH_3COOH \rightleftharpoons CH_3COOH_2^+ + Cl^-} \quad (\text{部分电离})
\]

在冰醋酸中，HClO\textsubscript{4} 仍是强酸，但 HCl 变成了弱酸---溶剂''区分''了它们。

\begin{quote}
\textbf{竞赛考点}：判断某溶剂是否有拉平效应，就看该溶剂能否接受质子。水是质子受体，所以有拉平效应。己烷不能接受质子，没有拉平效应。
\end{quote}

\begin{center}\rule{0.5\linewidth}{0.5pt}\end{center}

\subsection{三、共轭酸碱对与 Ka/Kb}\label{ux4e09ux5171ux8f6dux9178ux78b1ux5bf9ux4e0e-kakb}

\subsubsection{3.1 Ka 与 Kb 的关系}\label{ka-ux4e0e-kb-ux7684ux5173ux7cfb}

对共轭酸碱对 HA/A\textsuperscript{-}：

\[
K_a = \frac{[\mathrm{H_3O^+}][\mathrm{A^-}]}{[\mathrm{HA}]} \qquad K_b = \frac{[\mathrm{HA}][\mathrm{OH^-}]}{[\mathrm{A^-}]}
\]

两式相乘：

\[
\boxed{K_a \cdot K_b = K_w = 1.0 \times 10^{-14} \quad (25°C)}
\]

\[
\boxed{\mathrm{p}K_a + \mathrm{p}K_b = 14.00}
\]

\begin{quote}
\textbf{推导}：Ka · Kb = {[}H\textsubscript{3}O\textsuperscript{+}{]}{[}A\textsuperscript{-}{]}/{[}HA{]} × {[}HA{]}{[}OH\textsuperscript{-}{]}/{[}A\textsuperscript{-}{]} = {[}H\textsubscript{3}O\textsuperscript{+}{]}{[}OH\textsuperscript{-}{]} = Kw
\end{quote}

\subsubsection{3.2 Ka 的物理意义}\label{ka-ux7684ux7269ux7406ux610fux4e49}

\begin{itemize}
\tightlist
\item
  Ka 越大 \箭 酸越强 \箭 越容易给出质子
\item
  Ka 越小 \箭 酸越弱 \箭 越不容易给出质子
\item
  Ka \textgreater{} 1：强酸(在水中完全电离)
\item
  Ka \textless{} 10\textsuperscript{-3}：弱酸
\item
  Ka \textless{} 10\textsuperscript{-7}：极弱酸
\end{itemize}

\textbf{常见酸的 pKa 值}(25°C)：

{\def\LTcaptype{none} % do not increment counter
\begin{longtable}[]{@{}lcl@{}}
\toprule\noalign{}
酸 & pKa & 强度等级 \\
\midrule\noalign{}
\endhead
\bottomrule\noalign{}
\endlastfoot
HClO\textsubscript{4} & −10 & 极强酸 \\
HCl & −7 & 强酸 \\
HNO\textsubscript{3} & −1.4 & 强酸 \\
HF & 3.2 & 弱酸 \\
HAc & 4.76 & 弱酸 \\
H\textsubscript{2}CO\textsubscript{3} & 6.35(Ka\textsubscript{1}) & 极弱酸 \\
H\textsubscript{2}S & 7.02(Ka\textsubscript{1}) & 极弱酸 \\
HCN & 9.21 & 极弱酸 \\
H\textsubscript{2}O & 15.7 & 极极弱酸 \\
NH\textsubscript{3} & 38(作为酸) & 极极弱酸 \\
\end{longtable}
}

\begin{quote}
\textbf{易错点}：有机化学中常以 H\textsubscript{2}O 的 pKa = 15.7 作为参考，而无机化学中以 Kw = 10\textsuperscript{-14}(即 pKa = 14)作为参考。两者基准不同！
\end{quote}

\subsubsection{3.3 弱酸弱碱分类}\label{ux5f31ux9178ux5f31ux78b1ux5206ux7c7b}

{\def\LTcaptype{none} % do not increment counter
\begin{longtable}[]{@{}lll@{}}
\toprule\noalign{}
类型 & 示例 & 特点 \\
\midrule\noalign{}
\endhead
\bottomrule\noalign{}
\endlastfoot
\textbf{分子酸} & HF, HAc, H\textsubscript{2}S, HCN & 中性分子作为酸 \\
\textbf{阳离子酸} & NH\textsubscript{4}\textsuperscript{+}, {[}Al(H\textsubscript{2}O)\textsubscript{6}{]}\textsuperscript{3+} & 阳离子作为酸(水解) \\
\textbf{阴离子碱} & Ac\textsuperscript{-}, CN\textsuperscript{-}, S\textsuperscript{2-}, PO\textsubscript{4}\textsuperscript{3-} & 阴离子作为碱 \\
\textbf{两性物质} & H\textsubscript{2}O, HCO\textsubscript{3}\textsuperscript{-}, HPO\textsubscript{4}\textsuperscript{2-} & 既可作酸又可作碱 \\
\end{longtable}
}

\begin{quote}
\textbf{金属离子水解}：高价金属离子如 Al\textsuperscript{3+}、Fe\textsuperscript{3+} 在水中以水合离子形式存在，如 {[}Al(H\textsubscript{2}O)\textsubscript{6}{]}\textsuperscript{3+}。这些水合离子可以逐步电离出 H\textsuperscript{+}：

{[}Al(H\textsubscript{2}O)\textsubscript{6}{]}\textsuperscript{3+} \衡 [Al(H~2~O)~5~(OH)]\textsuperscript{2+} + H\textsuperscript{+} \衡 [Al(H~2~O)~4~(OH)~2~]\textsuperscript{+} + 2H\textsuperscript{+} \衡 \ldots{}

这就是为什么 AlCl\textsubscript{3} 溶液呈酸性。离子势 z/r(电荷/半径)越大，水解程度越大，酸性越强。
\end{quote}

\begin{center}\rule{0.5\linewidth}{0.5pt}\end{center}

\subsection{四、水的自偶电离与 pH}\label{ux56dbux6c34ux7684ux81eaux5076ux7535ux79bbux4e0e-ph}

\subsubsection{4.1 水的离子积}\label{ux6c34ux7684ux79bbux5b50ux79ef}

\[
\mathrm{H_2O + H_2O \rightleftharpoons H_3O^+ + OH^-}
\]

\[
\boxed{K_w = [\mathrm{H_3O^+}][\mathrm{OH^-}] = 1.0 \times 10^{-14} \quad (25°C)}
\]

\textbf{Kw 的温度依赖}：

{\def\LTcaptype{none} % do not increment counter
\begin{longtable}[]{@{}cc@{}}
\toprule\noalign{}
温度(°C) & Kw \\
\midrule\noalign{}
\endhead
\bottomrule\noalign{}
\endlastfoot
0 & 0.114 × 10\textsuperscript{-14} \\
10 & 0.293 × 10\textsuperscript{-14} \\
20 & 0.681 × 10\textsuperscript{-14} \\
25 & 1.00 × 10\textsuperscript{-14} \\
37 & 2.42 × 10\textsuperscript{-14} \\
50 & 5.47 × 10\textsuperscript{-14} \\
100 & 51.3 × 10\textsuperscript{-14} \\
\end{longtable}
}

\begin{quote}
\textbf{温度升高 \箭 Kw 增大}：水的自偶电离是吸热过程(ΔH \textgreater{} 0)，升温促进电离。
\textbf{37°C(体温)时 Kw = 2.42 × 10\textsuperscript{-14}}，此时中性 pH \约等于 6.81(不是 7.00！)。
\end{quote}

\subsubsection{4.2 pH 定义与有效数字}\label{ph-ux5b9aux4e49ux4e0eux6709ux6548ux6570ux5b57}

\[
\mathrm{pH} = -\lg[\mathrm{H_3O^+}]
\]

\textbf{有效数字规则}：pH 的有效数字位数由 {[}H\textsubscript{3}O\textsuperscript{+}{]} 的有效数字位数决定。

{\def\LTcaptype{none} % do not increment counter
\begin{longtable}[]{@{}ccc@{}}
\toprule\noalign{}
{[}H\textsubscript{3}O\textsuperscript{+}{]} & 有效数字 & pH \\
\midrule\noalign{}
\endhead
\bottomrule\noalign{}
\endlastfoot
1.0 × 10\textsuperscript{-3} & 2 位 & 3.00 \\
1 × 10\textsuperscript{-3} & 1 位 & 3 \\
0.00100 & 3 位 & 3.000 \\
\end{longtable}
}

\begin{quote}
\textbf{常见错误}：pH = 3.00 有 2 位有效数字，不是 3 位！小数点前的数字不计入有效数字。
\end{quote}

\begin{center}\rule{0.5\linewidth}{0.5pt}\end{center}

\subsection{五、弱酸弱碱电离平衡}\label{ux4e94ux5f31ux9178ux5f31ux78b1ux7535ux79bbux5e73ux8861}

\subsubsection{5.1 单质子弱酸：假设-验证法}\label{ux5355ux8d28ux5b50ux5f31ux9178ux5047ux8bbe-ux9a8cux8bc1ux6cd5}

\textbf{标准五步法}：

\begin{enumerate}
\def\labelenumi{\arabic{enumi}.}
\tightlist
\item
  写出电离方程式
\item
  列 ICE 表
\item
  写 Ka 表达式
\item
  假设求解(简化法 or 精确法)
\item
  验证假设合理性
\end{enumerate}

\textbf{简化公式}：

\[
[\mathrm{H_3O^+}] \approx \sqrt{K_a \cdot c}
\]

\begin{quote}
\textbf{适用条件}：\(c/K_a > 400\)(由 α \textless{} 5\% 推导：\(\frac{1-0.05}{0.05^2} = 380 \approx 400\))

同时需要 \(c \cdot K_a > 10 K_w\)(确保水的电离可忽略)。

\textbf{易错点}：如果 Ka 较大(如二氯乙酸 Ka = 4.5 × 10\textsuperscript{-2})，即使 c/Ka \textgreater{} 400 也可能不满足 α \textless{} 5\%，此时必须用精确法(二次方程)。
\end{quote}

\textbf{弱碱的对应公式}：

\[
[\mathrm{OH^-}] \approx \sqrt{K_b \cdot c} \qquad (c/K_b > 400)
\]

\subsubsection{5.2 精确法(二次方程)}\label{ux7cbeux786eux6cd5ux4e8cux6b21ux65b9ux7a0b}

当简化条件不满足时，解精确方程：

\[
K_a = \frac{x^2}{c - x}
\]

\[
x^2 + K_a x - K_a c = 0
\]

\[
x = \frac{-K_a + \sqrt{K_a^2 + 4K_a c}}{2}
\]

\subsubsection{5.3 反例：简化公式失效}\label{ux53cdux4f8bux7b80ux5316ux516cux5f0fux5931ux6548}

\textbf{Na\textsubscript{2}S 溶液}：Kb\textsubscript{1}(S\textsuperscript{2-}) = 8.3 × 10\textsuperscript{-2}，c = 0.10 M

c/Kb\textsubscript{1} = 0.10/0.083 = 1.2 ≪ 400

简化公式给出 {[}OH\textsuperscript{-}{]} = √(0.083 × 0.10) = 0.091 M(误差很大！)

精确法：x\textsuperscript{2}/(0.10 − x) = 0.083 \箭 x\textsuperscript{2} + 0.083x − 0.0083 = 0 \箭 x = 0.059 M

\begin{quote}
\textbf{教训}：使用简化公式前，\textbf{必须}先检查 c/Ka \textgreater{} 400(或 c/Kb \textgreater{} 400)这个条件。``每次用简化公式前，先检查两个条件''是酸碱计算的第一条铁律。
\end{quote}

\subsubsection{\texorpdfstring{5.4 弱酸溶液中 {[}H\textsuperscript{+}{]} 的精确处理}{5.4 弱酸溶液中 {[}H+{]} 的精确处理}}\label{ux5f31ux9178ux6eb6ux6db2ux4e2d-h-ux7684ux7cbeux786eux5904ux7406}

对弱酸 HA 溶液，精确的质子平衡方程(PBE)为：

\[
[\mathrm{H^+}] = [\mathrm{A^-}] + [\mathrm{OH^-}]
\]

即 H\textsuperscript{+} 既来自 HA 的电离，也来自水的电离。当 c·Ka \textgreater{} 10Kw 时，{[}OH\textsuperscript{-}{]} 可忽略，简化为 {[}H\textsuperscript{+}{]} \约等于 [A^-^]。

\begin{center}\rule{0.5\linewidth}{0.5pt}\end{center}

\subsection{六、多元弱酸与两性物质}\label{ux516dux591aux5143ux5f31ux9178ux4e0eux4e24ux6027ux7269ux8d28}

\subsubsection{6.1 多元弱酸的分级电离}\label{ux591aux5143ux5f31ux9178ux7684ux5206ux7ea7ux7535ux79bb}

以 H\textsubscript{2}S 为例：

\[
\mathrm{H_2S \rightleftharpoons H^+ + HS^-} \quad K_{a1} = 1.0 \times 10^{-7}
\]
\[
\mathrm{HS^- \rightleftharpoons H^+ + S^{2-}} \quad K_{a2} = 1.1 \times 10^{-12}
\]

\textbf{关键规则}：Ka\textsubscript{1} ≫ Ka\textsubscript{2}(通常差 10\textsuperscript{3}\textasciitilde10\textsuperscript{5} 倍)，因此：
- {[}H\textsuperscript{+}{]} 主要由第一步电离决定
- {[}S\textsuperscript{2-}{]} \约等于 Ka\textsubscript{2}(与 c 无关！)

\[
[\mathrm{H^+}] \approx \sqrt{K_{a1} \cdot c} \qquad [\mathrm{S^{2-}}] \approx K_{a2}
\]

\textbf{H\textsubscript{2}S 物种分布与 pH 的关系}：

{\def\LTcaptype{none} % do not increment counter
\begin{longtable}[]{@{}ll@{}}
\toprule\noalign{}
pH 范围 & 主导物种 \\
\midrule\noalign{}
\endhead
\bottomrule\noalign{}
\endlastfoot
\textless{} 7 & H\textsubscript{2}S \\
7 \textasciitilde{} 12 & HS\textsuperscript{-} \\
\textgreater{} 12 & S\textsuperscript{2-} \\
\end{longtable}
}

\subsubsection{6.2 两性物质的 pH}\label{ux4e24ux6027ux7269ux8d28ux7684-ph}

对 NaHCO\textsubscript{3} 溶液，HCO\textsubscript{3}\textsuperscript{-} 既电离又水解：

\[
\mathrm{HCO_3^- \rightleftharpoons H^+ + CO_3^{2-}} \quad K_{a2}
\]
\[
\mathrm{HCO_3^- + H_2O \rightleftharpoons H_2CO_3 + OH^-} \quad K_b = K_w/K_{a1}
\]

\[
\boxed{[\mathrm{H^+}] \approx \sqrt{K_{a1} \cdot K_{a2}}}
\]

\[
\boxed{\mathrm{pH} = \frac{1}{2}(\mathrm{p}K_{a1} + \mathrm{p}K_{a2})}
\]

\textbf{实例}：0.20 M NaHCO\textsubscript{3}

pKa\textsubscript{1} = 6.35, pKa\textsubscript{2} = 10.33

pH = ½(6.35 + 10.33) = 8.34

\begin{quote}
\textbf{注意}：这个公式与浓度无关！只要 Ka\textsubscript{1} 和 Ka\textsubscript{2} 相差足够大(通常 \textgreater{} 10\textsuperscript{3})，且 c/Ka\textsubscript{1} \textgreater{} 400，公式就成立。

反例：NaHC\textsubscript{2}O\textsubscript{4}(草酸氢钠)的 Ka\textsubscript{2} = 5.4 × 10\textsuperscript{-5}，c/Ka\textsubscript{2} \textless{} 400 时简化公式失效，需要精确处理。
\end{quote}

\begin{center}\rule{0.5\linewidth}{0.5pt}\end{center}

\subsection{七、缓冲溶液}\label{ux4e03ux7f13ux51b2ux6eb6ux6db2}

\subsubsection{7.1 Henderson-Hasselbalch 方程}\label{henderson-hasselbalch-ux65b9ux7a0b}

对弱酸 HA 及其共轭碱 A\textsuperscript{-} 的混合溶液：

\[
\boxed{\mathrm{pH} = \mathrm{p}K_a + \lg\frac{[\mathrm{A^-}]}{[\mathrm{HA}]}}
\]

对弱碱 B 及其共轭酸 BH\textsuperscript{+} 的混合溶液：

\[
\mathrm{pOH} = \mathrm{p}K_b + \lg\frac{[\mathrm{BH^+}]}{[\mathrm{B}]}
\]

\subsubsection{7.2 缓冲溶液的设计原则}\label{ux7f13ux51b2ux6eb6ux6db2ux7684ux8bbeux8ba1ux539fux5219}

{\def\LTcaptype{none} % do not increment counter
\begin{longtable}[]{@{}ll@{}}
\toprule\noalign{}
原则 & 说明 \\
\midrule\noalign{}
\endhead
\bottomrule\noalign{}
\endlastfoot
\textbf{选对 pKa} & pKa 应接近目标 pH(pKa ± 1 范围内) \\
\textbf{比例接近 1:1} & {[}A\textsuperscript{-}{]}/{[}HA{]} = 1 时缓冲容量最大 \\
\textbf{浓度足够大} & 浓度越高，缓冲容量越大 \\
\end{longtable}
}

\textbf{常用缓冲溶液}：

{\def\LTcaptype{none} % do not increment counter
\begin{longtable}[]{@{}lcc@{}}
\toprule\noalign{}
缓冲体系 & pKa & 有效缓冲范围 \\
\midrule\noalign{}
\endhead
\bottomrule\noalign{}
\endlastfoot
HCOOH/HCOO\textsuperscript{-} & 3.75 & 2.75 \textasciitilde{} 4.75 \\
HAc/Ac\textsuperscript{-} & 4.76 & 3.76 \textasciitilde{} 5.76 \\
H\textsubscript{2}CO\textsubscript{3}/HCO\textsubscript{3}\textsuperscript{-} & 6.35 & 5.35 \textasciitilde{} 7.35 \\
H\textsubscript{2}PO\textsubscript{4}\textsuperscript{-}/HPO\textsubscript{4}\textsuperscript{2-} & 7.20 & 6.20 \textasciitilde{} 8.20 \\
NH\textsubscript{4}\textsuperscript{+}/NH\textsubscript{3} & 9.25 & 8.25 \textasciitilde{} 10.25 \\
HPO\textsubscript{4}\textsuperscript{2-}/PO\textsubscript{4}\textsuperscript{3-} & 12.32 & 11.32 \textasciitilde{} 13.32 \\
\end{longtable}
}

\subsubsection{7.3 缓冲容量的定量分析}\label{ux7f13ux51b2ux5bb9ux91cfux7684ux5b9aux91cfux5206ux6790}

缓冲容量 β 定义为使 pH 改变 1 个单位所需加入的强酸或强碱的量。

\textbf{关键结论}：
- 当 {[}A\textsuperscript{-}{]}/{[}HA{]} = 1(即 pH = pKa)时，缓冲容量最大
- 当 {[}A\textsuperscript{-}{]}/{[}HA{]} = 99:1 或 1:99 时，缓冲容量接近零
- 缓冲溶液的总浓度越高，缓冲容量越大

\textbf{实例}：向 1 L 缓冲溶液中加入 0.01 mol H\textsubscript{3}O\textsuperscript{+}

{\def\LTcaptype{none} % do not increment counter
\begin{longtable}[]{@{}cccc@{}}
\toprule\noalign{}
初始比例 {[}A\textsuperscript{-}{]}/{[}HA{]} & 初始 pH & 加入后 pH & ΔpH \\
\midrule\noalign{}
\endhead
\bottomrule\noalign{}
\endlastfoot
1:1 & 4.76 & 4.74 & 0.02 \\
10:1 & 5.76 & 5.53 & 0.23 \\
99:1 & 6.76 & 5.00 & 1.76 \\
\end{longtable}
}

\begin{quote}
\textbf{比例 1:1 时缓冲能力最强}，这是设计缓冲溶液的黄金法则。
\end{quote}

\subsubsection{7.4 缓冲溶液的实际应用}\label{ux7f13ux51b2ux6eb6ux6db2ux7684ux5b9eux9645ux5e94ux7528}

\textbf{血液缓冲系统}：人体血液 pH = 7.35 \textasciitilde{} 7.45，主要由 HCO\textsubscript{3}\textsuperscript{-}/H\textsubscript{2}CO\textsubscript{3} 缓冲体系维持。

\[
\mathrm{CO_2 + H_2O \rightleftharpoons H_2CO_3 \rightleftharpoons H^+ + HCO_3^-}
\]

\[
\mathrm{pH} = \mathrm{p}K_{a1} + \lg\frac{[\mathrm{HCO_3^-}]}{[\mathrm{H_2CO_3}]} = 6.10 + \lg\frac{24}{1.2} = 7.40
\]

(正常血液中 {[}HCO\textsubscript{3}\textsuperscript{-}{]} = 24 mmol/L，{[}H\textsubscript{2}CO\textsubscript{3}{]} = 1.2 mmol/L)

\textbf{血红蛋白的氧结合平衡}：

\[
\mathrm{HbH^+ + O_2 \rightleftharpoons HbO_2 + H^+}
\]

当 O\textsubscript{2} 与 Hb 结合时释放 H\textsuperscript{+}，H\textsuperscript{+} 又被 HCO\textsubscript{3}\textsuperscript{-} 缓冲---这就是 Bohr 效应的化学基础。

\begin{quote}
\textbf{缓冲的''海绵''类比}：缓冲溶液就像一块海绵---倒入一点酸，它吸收了，pH 几乎不变；倒入一点碱，它也吸收了，pH 也几乎不变。但如果倒入太多的酸或碱，海绵饱和了，pH 就会急剧变化。
\end{quote}

\begin{center}\rule{0.5\linewidth}{0.5pt}\end{center}

\subsection{八、Lewis 酸碱与 HSAB 原理}\label{ux516blewis-ux9178ux78b1ux4e0e-hsab-ux539fux7406}

\subsubsection{8.1 Lewis 酸碱}\label{lewis-ux9178ux78b1}

\textbf{Lewis 酸}：电子对受体(能接受电子对的物质)
\textbf{Lewis 碱}：电子对供体(能给出电子对的物质)

{\def\LTcaptype{none} % do not increment counter
\begin{longtable}[]{@{}lll@{}}
\toprule\noalign{}
Lewis 酸 & Lewis 碱 & 反应 \\
\midrule\noalign{}
\endhead
\bottomrule\noalign{}
\endlastfoot
BF\textsubscript{3} & NH\textsubscript{3} & F\textsubscript{3}B\左箭NH\textsubscript{3}(配位键) \\
H\textsuperscript{+} & H\textsubscript{2}O & H\textsubscript{3}O\textsuperscript{+} \\
Cu\textsuperscript{2+} & NH\textsubscript{3} & {[}Cu(NH\textsubscript{3})\textsubscript{4}{]}\textsuperscript{2+} \\
AlCl\textsubscript{3} & Cl\textsuperscript{-} & AlCl\textsubscript{4}\textsuperscript{-} \\
\end{longtable}
}

\begin{quote}
Lewis 酸碱理论是最广义的酸碱定义，涵盖了所有 Brønsted 酸碱反应(因为 H\textsuperscript{+} 本身就是 Lewis 酸)。
\end{quote}

\subsubsection{8.2 HSAB 软硬分类}\label{hsab-ux8f6fux786cux5206ux7c7b}

\textbf{Hard and Soft Acids and Bases(软硬酸碱理论)}：

{\def\LTcaptype{none} % do not increment counter
\begin{longtable}[]{@{}
  >{\raggedright\arraybackslash}p{(\linewidth - 4\tabcolsep) * \real{0.3333}}
  >{\raggedright\arraybackslash}p{(\linewidth - 4\tabcolsep) * \real{0.3333}}
  >{\raggedright\arraybackslash}p{(\linewidth - 4\tabcolsep) * \real{0.3333}}@{}}
\toprule\noalign{}
\begin{minipage}[b]{\linewidth}\raggedright
分类
\end{minipage} & \begin{minipage}[b]{\linewidth}\raggedright
酸
\end{minipage} & \begin{minipage}[b]{\linewidth}\raggedright
碱
\end{minipage} \\
\midrule\noalign{}
\endhead
\bottomrule\noalign{}
\endlastfoot
\textbf{硬酸} & H\textsuperscript{+}, Li\textsuperscript{+}, Na\textsuperscript{+}, Mg\textsuperscript{2+}, Al\textsuperscript{3+}, Fe\textsuperscript{3+}, BF\textsubscript{3} & H\textsubscript{2}O, OH\textsuperscript{-}, F\textsuperscript{-}, Cl\textsuperscript{-}, NH\textsubscript{3}, ROH \\
\textbf{交界酸} & Fe\textsuperscript{2+}, Co\textsuperscript{2+}, Ni\textsuperscript{2+}, Cu\textsuperscript{2+}, Zn\textsuperscript{2+} & Br\textsuperscript{-}, N\textsubscript{3}\textsuperscript{-}, NO\textsubscript{2}\textsuperscript{-}, SO\textsubscript{3}\textsuperscript{2-} \\
\textbf{软酸} & Cu\textsuperscript{+}, Ag\textsuperscript{+}, Au\textsuperscript{+}, Hg\textsuperscript{2+}, Pd\textsuperscript{2+}, Pt\textsuperscript{2+} & H\textsuperscript{-}, I\textsuperscript{-}, CN\textsuperscript{-}, CO, S\textsuperscript{2-}, RS\textsuperscript{-} \\
\end{longtable}
}

\textbf{HSAB 亲和规则}：

\begin{quote}
\textbf{硬酸倾向与硬碱结合，软酸倾向与软碱结合。}
\end{quote}

{\def\LTcaptype{none} % do not increment counter
\begin{longtable}[]{@{}lll@{}}
\toprule\noalign{}
组合 & 稳定性 & 示例 \\
\midrule\noalign{}
\endhead
\bottomrule\noalign{}
\endlastfoot
硬酸 + 硬碱 & \textbf{稳定} & NaF(硬酸 Na\textsuperscript{+} + 硬碱 F\textsuperscript{-}) \\
软酸 + 软碱 & \textbf{稳定} & AgI(软酸 Ag\textsuperscript{+} + 软碱 I\textsuperscript{-}) \\
硬酸 + 软碱 & 不稳定 & AgF(软酸 Ag\textsuperscript{+} + 硬碱 F\textsuperscript{-}) \\
软酸 + 硬碱 & 不稳定 & NaI(硬酸 Na\textsuperscript{+} + 软碱 I\textsuperscript{-}) \\
\end{longtable}
}

\begin{quote}
\textbf{HSAB 的''媒人''类比}：硬酸找硬碱就像两个''固执''的人在一起---稳定。软酸找软碱就像两个''温柔''的人在一起---也稳定。但硬配软就像性格不合，容易''分手''。

\textbf{竞赛应用}：为什么 AgI 不溶于水而 AgF 溶于水？因为 Ag\textsuperscript{+} 是软酸，I\textsuperscript{-} 是软碱(稳定结合)，而 F\textsuperscript{-} 是硬碱(不稳定结合)。
\end{quote}

\begin{center}\rule{0.5\linewidth}{0.5pt}\end{center}

\subsection{九、中和反应与电离常数}\label{ux4e5dux4e2dux548cux53cdux5e94ux4e0eux7535ux79bbux5e38ux6570}

\subsubsection{9.1 中和反应的平衡常数}\label{ux4e2dux548cux53cdux5e94ux7684ux5e73ux8861ux5e38ux6570}

{\def\LTcaptype{none} % do not increment counter
\begin{longtable}[]{@{}llcl@{}}
\toprule\noalign{}
反应类型 & 示例 & K 值 & 说明 \\
\midrule\noalign{}
\endhead
\bottomrule\noalign{}
\endlastfoot
强酸 + 强碱 & HCl + NaOH & 1/Kw = 10\textsuperscript{14} & 反应极其彻底 \\
强酸 + 弱碱 & HCl + NH\textsubscript{3} & Ka/Kw = 10\textsuperscript{9}·\textsuperscript{25} & 反应较彻底 \\
弱酸 + 强碱 & HAc + NaOH & Kb/Kw = 10\textsuperscript{9}·\textsuperscript{24} & 反应较彻底 \\
弱酸 + 弱碱 & HAc + NH\textsubscript{3} & Ka·Kb/Kw\textsuperscript{2} = 10\textsuperscript{4}·\textsuperscript{76} & 反应程度有限 \\
\end{longtable}
}

\begin{quote}
\textbf{规律}：反应物酸碱越强，中和反应的 K 越大，反应越彻底。
\end{quote}

\subsubsection{9.2 弱-弱中和的实例}\label{ux5f31-ux5f31ux4e2dux548cux7684ux5b9eux4f8b}

\[
\mathrm{H_2CO_3 + SiO_3^{2-} \to HCO_3^- + HSiO_3^-}
\]

K = Ka\textsubscript{1}(H\textsubscript{2}CO\textsubscript{3})/Ka\textsubscript{2}(H\textsubscript{2}SiO\textsubscript{3}) \约等于 10\textsuperscript{-6}·\textsuperscript{35}/10\textsuperscript{-12}·\textsuperscript{8} = 10\textsuperscript{6}·\textsuperscript{45}

反应较彻底(K \textgreater{} 10\textsuperscript{6})，这就是为什么 CO\textsubscript{2} 通入硅酸钠溶液会产生硅酸沉淀。

\begin{center}\rule{0.5\linewidth}{0.5pt}\end{center}

\subsection{十、知识网络与方法论}\label{ux5341ux77e5ux8bc6ux7f51ux7edcux4e0eux65b9ux6cd5ux8bba}

\subsubsection{10.1 弱酸弱碱计算决策流程}\label{ux5f31ux9178ux5f31ux78b1ux8ba1ux7b97ux51b3ux7b56ux6d41ux7a0b}

\begin{verbatim}
已知 Ka 和 c
  → 检查 c/Ka > 400？
    → 是：用简化公式 [H⁺] = √(Ka·c)
    → 否：用精确法（二次方程）
  → 检查 c·Ka > 10Kw？
    → 是：忽略水的电离
    → 否：需考虑水的自偶电离
\end{verbatim}

\subsubsection{10.2 两性物质计算决策}\label{ux4e24ux6027ux7269ux8d28ux8ba1ux7b97ux51b3ux7b56}

\begin{verbatim}
物质是两性物质（如 NaHCO₃）
  → 检查 Ka₂ 和 Kb = Kw/Ka₁ 的大小关系
  → 若 Ka₂ ≈ Kb：用 [H⁺] = √(Ka₁·Ka₂)
  → 若 Ka₂ ≫ Kb：视为弱酸处理
  → 若 Kb ≫ Ka₂：视为弱碱处理
\end{verbatim}

\subsubsection{10.3 酸碱计算的''五个问题''}\label{ux9178ux78b1ux8ba1ux7b97ux7684ux4e94ux4e2aux95eeux9898}

在做任何酸碱计算前，先问自己：

\begin{enumerate}
\def\labelenumi{\arabic{enumi}.}
\tightlist
\item
  \textbf{溶液中有哪些物种？}(分子、离子、两性物质)
\item
  \textbf{发生了哪些反应？}(电离、水解、中和)
\item
  \textbf{哪些可以忽略？}(水的电离、第二步电离等)
\item
  \textbf{哪些假设是合理的？}(c/Ka \textgreater{} 400、α \textless{} 5\%)
\item
  \textbf{答案是否合理？}(pH 范围、浓度正负)
\end{enumerate}

\begin{center}\rule{0.5\linewidth}{0.5pt}\end{center}

\subsection{十一、典型例题}\label{ux5341ux4e00ux5178ux578bux4f8bux9898}

\textbf{题干}：指出下列反应中的共轭酸碱对：

\[
\mathrm{H_2S + CN^- \rightleftharpoons HS^- + HCN}
\]

\textbf{解析}：

{\def\LTcaptype{none} % do not increment counter
\begin{longtable}[]{@{}lll@{}}
\toprule\noalign{}
反应物 & 角色 & 共轭对 \\
\midrule\noalign{}
\endhead
\bottomrule\noalign{}
\endlastfoot
H\textsubscript{2}S \箭 HS\textsuperscript{-} & 酸 \箭 共轭碱 & H\textsubscript{2}S / HS\textsuperscript{-} \\
CN\textsuperscript{-} \箭 HCN & 碱 \箭 共轭酸 & HCN / CN\textsuperscript{-} \\
\end{longtable}
}

\textbf{答案}：两对共轭酸碱对为 (H\textsubscript{2}S/HS\textsuperscript{-}) 和 (HCN/CN\textsuperscript{-})。

\begin{center}\rule{0.5\linewidth}{0.5pt}\end{center}

\textbf{题干}：计算 0.10 M HF 溶液的 pH(Ka = 7.2 × 10\textsuperscript{-4})。

\textbf{解析}：

c/Ka = 0.10/7.2×10\textsuperscript{-4} = 139 \textless{} 400，简化公式不适用，用精确法。

设 {[}H\textsuperscript{+}{]} = x：

\[K_a = \frac{x^2}{0.10 - x} = 7.2 \times 10^{-4}\]

\[x^2 + 7.2 \times 10^{-4} x - 7.2 \times 10^{-5} = 0\]

\[x = \frac{-7.2 \times 10^{-4} + \sqrt{(7.2 \times 10^{-4})^2 + 4 \times 7.2 \times 10^{-5}}}{2} = 7.8 \times 10^{-3}\]

\[\mathrm{pH} = -\lg(7.8 \times 10^{-3}) = 2.11\]

\[\alpha = \frac{7.8 \times 10^{-3}}{0.10} \times 100\% = 7.8\%\]

\textbf{答案}：pH = 2.11，α = 7.8\%。

\begin{quote}
如果用简化公式：{[}H\textsuperscript{+}{]} = √(7.2×10\textsuperscript{-4} × 0.10) = 8.5×10\textsuperscript{-3}，pH = 2.07，误差约 9\%。c/Ka = 139 \textless{} 400 时误差不可忽略。
\end{quote}

\begin{center}\rule{0.5\linewidth}{0.5pt}\end{center}

\textbf{题干}：计算 0.20 M NaHCO\textsubscript{3} 溶液的 pH(Ka\textsubscript{1} = 4.47 × 10\textsuperscript{-7}，Ka\textsubscript{2} = 4.68 × 10\textsuperscript{-11})。

\textbf{解析}：

\[\mathrm{pH} = \frac{1}{2}(\mathrm{p}K_{a1} + \mathrm{p}K_{a2}) = \frac{1}{2}(6.35 + 10.33) = 8.34\]

\textbf{答案}：pH = 8.34(碱性，符合 HCO\textsubscript{3}\textsuperscript{-} 水解大于电离的预期)。

\begin{center}\rule{0.5\linewidth}{0.5pt}\end{center}

\textbf{题干}：用 HAc 和 NaAc 配制 pH = 5.00 的缓冲溶液 1.00 L，需要多少克 NaAc 和多少毫升冰醋酸(密度 1.05 g/mL)？(Ka = 1.75 × 10\textsuperscript{-5})

\textbf{解析}：

\[\lg\frac{[\mathrm{Ac^-}]}{[\mathrm{HAc}]} = \mathrm{pH} - \mathrm{p}K_a = 5.00 - 4.76 = 0.24\]

\[\frac{[\mathrm{Ac^-}]}{[\mathrm{HAc}]} = 10^{0.24} = 1.74\]

设 {[}HAc{]} = x，则 {[}Ac\textsuperscript{-}{]} = 1.74x，且 x + 1.74x = c\_total。

取 c\_total = 0.50 M(任意选择，影响缓冲容量)：

x = 0.18 M(HAc)，1.74x = 0.31 M(Ac\textsuperscript{-})

NaAc 质量 = 0.31 × 1.00 × 82.03 = 25.4 g
冰醋酸体积 = 0.18 × 1.00 × 60.05 / 1.05 = 10.3 mL

\textbf{答案}：需要 25.4 g NaAc 和 10.3 mL 冰醋酸。

\begin{center}\rule{0.5\linewidth}{0.5pt}\end{center}

\textbf{题干}：用 HSAB 原理预测以下配离子的稳定性顺序：
(a) {[}AgF\textsubscript{2}{]}\textsuperscript{-} vs {[}AgI\textsubscript{2}{]}\textsuperscript{-}
(b) {[}AlF\textsubscript{6}{]}\textsuperscript{3-} vs {[}AlI\textsubscript{6}{]}\textsuperscript{3-}

\textbf{解析}：

\begin{enumerate}
\def\labelenumi{(\alph{enumi})}
\item
  Ag\textsuperscript{+} 是软酸，F\textsuperscript{-} 是硬碱，I\textsuperscript{-} 是软碱。
  {[}AgI\textsubscript{2}{]}\textsuperscript{-}(软-软)比 {[}AgF\textsubscript{2}{]}\textsuperscript{-}(软-硬)稳定得多。
\item
  Al\textsuperscript{3+} 是硬酸，F\textsuperscript{-} 是硬碱，I\textsuperscript{-} 是软碱。
  {[}AlF\textsubscript{6}{]}\textsuperscript{3-}(硬-硬)比 {[}AlI\textsubscript{6}{]}\textsuperscript{3-}(硬-软)稳定得多。
\end{enumerate}

\textbf{答案}：(a) {[}AgI\textsubscript{2}{]}\textsuperscript{-} \textgreater{} {[}AgF\textsubscript{2}{]}\textsuperscript{-}；(b) {[}AlF\textsubscript{6}{]}\textsuperscript{3-} \textgreater{} {[}AlI\textsubscript{6}{]}\textsuperscript{3-}。

\begin{center}\rule{0.5\linewidth}{0.5pt}\end{center}

\subsection{十二、本节总结速查}\label{ux5341ux4e8cux672cux8282ux603bux7ed3ux901fux67e5}

{\def\LTcaptype{none} % do not increment counter
\begin{longtable}[]{@{}
  >{\raggedright\arraybackslash}p{(\linewidth - 4\tabcolsep) * \real{0.3333}}
  >{\raggedright\arraybackslash}p{(\linewidth - 4\tabcolsep) * \real{0.3333}}
  >{\raggedright\arraybackslash}p{(\linewidth - 4\tabcolsep) * \real{0.3333}}@{}}
\toprule\noalign{}
\begin{minipage}[b]{\linewidth}\raggedright
核心概念
\end{minipage} & \begin{minipage}[b]{\linewidth}\raggedright
关键公式
\end{minipage} & \begin{minipage}[b]{\linewidth}\raggedright
注意点
\end{minipage} \\
\midrule\noalign{}
\endhead
\bottomrule\noalign{}
\endlastfoot
共轭酸碱对 & Ka · Kb = Kw & pKa + pKb = 14 \\
弱酸 pH & {[}H\textsuperscript{+}{]} = √(Ka·c) & 条件：c/Ka \textgreater{} 400 且 c·Ka \textgreater{} 10Kw \\
两性物质 & {[}H\textsuperscript{+}{]} = √(Ka\textsubscript{1}·Ka\textsubscript{2}) & 条件：c/Ka\textsubscript{1} \textgreater{} 400 且 Ka\textsubscript{1}/Ka\textsubscript{2} \textgreater{} 10\textsuperscript{3} \\
缓冲溶液 & pH = pKa + lg({[}A\textsuperscript{-}{]}/{[}HA{]}) & pKa ± 1 范围有效 \\
Kw 温度依赖 & Kw 随温度升高而增大 & 37°C 时 Kw = 2.42×10\textsuperscript{-14} \\
拉平效应 & 溶剂将强酸''拉平''到同一水平 & 水对强酸有拉平效应 \\
HSAB & 硬亲硬，软亲软 & AgI 不溶于水，AgF 溶于水 \\
中和 K & 强酸+强碱 K=10\textsuperscript{14} & 弱酸+弱碱 K 可能很小 \\
\end{longtable}
}

\begin{center}\rule{0.5\linewidth}{0.5pt}\end{center}

\subsection{十三、三级练习题}\label{ux5341ux4e09ux4e09ux7ea7ux7ec3ux4e60ux9898}

\textbf{1.} 写出 NH\textsubscript{3} 在水中的 Brønsted 酸碱反应，指出共轭酸碱对。

\textbf{2.} 已知 Ka(HAc) = 1.75 × 10\textsuperscript{-5}，求 Kb(Ac\textsuperscript{-})。

\textbf{3.} 判断下列说法的正误：
(a) 强酸的共轭碱也是强碱。
(b) 缓冲溶液的 pH 恒定不变。
(c) pH = 7 的溶液一定是中性溶液。
(d) 两性物质溶液的 pH 与浓度无关。

\textbf{4.} 计算 0.10 M NH\textsubscript{3} 溶液的 pH(Kb = 1.77 × 10\textsuperscript{-5})。

\textbf{5.} 计算 0.20 M Na\textsubscript{2}CO\textsubscript{3} 溶液的 pH(Ka\textsubscript{1} = 4.47 × 10\textsuperscript{-7}，Ka\textsubscript{2} = 4.68 × 10\textsuperscript{-11})。

\textbf{6.} 用 HAc 和 NaAc 配制 pH = 5.00 的缓冲溶液，{[}Ac\textsuperscript{-}{]}/{[}HAc{]} = ?

\textbf{7.} 什么是拉平效应？为什么 HClO\textsubscript{4}、HCl、HNO\textsubscript{3} 在水中都是强酸？

\textbf{8.} 根据 HSAB 理论，预测 NaF 和 NaI 哪个更稳定。

\textbf{9.} 0.10 M HF 溶液中，c/Ka = 139 \textless{} 400。说明为什么不能用简化公式，并给出正确的处理方法。

\textbf{10.} 人体血液 pH = 7.40，主要由 HCO\textsubscript{3}\textsuperscript{-}/H\textsubscript{2}CO\textsubscript{3} 缓冲。若 {[}HCO\textsubscript{3}\textsuperscript{-}{]} = 24 mmol/L，{[}H\textsubscript{2}CO\textsubscript{3}{]} = 1.2 mmol/L，验证 pH。

\begin{center}\rule{0.5\linewidth}{0.5pt}\end{center}

\textbf{11.} \textbf{(混合弱酸)} 计算 0.10 M HAc 与 0.10 M HF 混合溶液的 pH(Ka(HAc) = 1.75×10\textsuperscript{-5}，Ka(HF) = 7.2×10\textsuperscript{-4})。

\textbf{12.} \textbf{(共离子效应)} 在 0.10 M HAc 中加入 NaAc 固体使 {[}NaAc{]} = 1.0 M，计算新 pH 和 HAc 的电离度。

\textbf{13.} \textbf{(缓冲容量)} 向 1 L pH = 4.76 的 HAc/NaAc 缓冲溶液({[}Ac\textsuperscript{-}{]}/{[}HAc{]} = 1:1)中加入 0.01 mol HCl，计算新 pH。

\textbf{14.} \textbf{(多元酸分步电离)} 0.10 M H\textsubscript{3}PO\textsubscript{4} 溶液中，计算 {[}H\textsuperscript{+}{]}、{[}H\textsubscript{2}PO\textsubscript{4}\textsuperscript{-}{]}、{[}HPO\textsubscript{4}\textsuperscript{2-}{]}、{[}PO\textsubscript{4}\textsuperscript{3-}{]}。已知 Ka\textsubscript{1}=7.1×10\textsuperscript{-3}，Ka\textsubscript{2}=6.3×10\textsuperscript{-8}，Ka\textsubscript{3}=4.2×10\textsuperscript{-13}。

\textbf{15.} \textbf{(弱酸电离度)} 某弱酸 HA 的 0.10 M 溶液 pH = 3.00，求 Ka 和 α。

\textbf{16.} \textbf{(缓冲溶液配制)} 需要配制 pH = 9.25 的缓冲溶液 500 mL，选择 NH\textsubscript{3}/NH\textsubscript{4}Cl 体系。计算所需 NH\textsubscript{4}Cl 的质量和浓氨水(15 M)的体积。

\textbf{17.} \textbf{(Na\textsubscript{2}S 精确计算)} 0.10 M Na\textsubscript{2}S 溶液，Kb\textsubscript{1} = 8.3×10\textsuperscript{-2}。用精确法计算 {[}OH\textsuperscript{-}{]} 和 pH。

\textbf{18.} \textbf{(拉平效应分析)} 在液氨溶剂中，HCl 和 HAc 都是弱酸。解释为什么。

\textbf{19.} \textbf{(非水溶剂酸碱)} 在液氨中，NaNH\textsubscript{2} 是强碱，NH\textsubscript{4}Cl 是强酸。写出它们在液氨中的酸碱反应。

\begin{center}\rule{0.5\linewidth}{0.5pt}\end{center}

\textbf{20.} \textbf{(混合缓冲体系)} 用 0.10 M H\textsubscript{3}PO\textsubscript{4} 和 0.10 M Na\textsubscript{3}PO\textsubscript{4} 配制 pH = 7.21 的缓冲溶液。需要取多少体积的 H\textsubscript{3}PO\textsubscript{4} 和多少体积的 Na\textsubscript{3}PO\textsubscript{4}？(Ka\textsubscript{1}=7.1×10\textsuperscript{-3}，Ka\textsubscript{2}=6.3×10\textsuperscript{-8}，Ka\textsubscript{3}=4.2×10\textsuperscript{-13})

\textbf{21.} \textbf{(多元酸完整物种计算)} 计算 0.10 M H\textsubscript{3}PO\textsubscript{4} 溶液中所有含磷物种的浓度，并绘制分布系数图(δ vs pH)。

\textbf{22.} \textbf{(血红蛋白缓冲)} 正常动脉血 pH = 7.40，{[}HCO\textsubscript{3}\textsuperscript{-}{]} = 24 mmol/L。已知 CO\textsubscript{2} 在血液中的溶解度为 0.030 mmol/(L·mmHg)，pKa\textsubscript{1} = 6.10。计算血液中 CO\textsubscript{2} 的分压。

\textbf{23.} \textbf{(弱酸 Ka 测定)} 用 0.1000 M NaOH 滴定 25.00 mL 未知弱酸 HA。当加入 12.50 mL NaOH 时，pH = 4.00；当加入 25.00 mL NaOH 时，pH = 8.00。求 HA 的 Ka。

\textbf{24.} \textbf{(HSAB 与溶解性)} 解释以下现象：
(a) AgCl 溶于浓氨水，AgI 不溶于浓氨水。
(b) HgS 不溶于浓 HNO\textsubscript{3} 和王水，但溶于浓 Na\textsubscript{2}S 溶液。
(c) HgCl\textsubscript{2} 是弱电解质，而 HgF\textsubscript{2} 是强电解质。

\textbf{25.} \textbf{(综合计算)} 25°C 时，向 100 mL 0.10 M HAc 溶液中逐滴加入 0.10 M NaOH。
(a) 计算未加 NaOH 时的 pH。
(b) 计算加入 50 mL NaOH 后的 pH(半中和点)。
(c) 计算加入 100 mL NaOH 后的 pH(等当点)。
(d) 计算加入 150 mL NaOH 后的 pH。
(e) 选择合适的指示剂，并解释为什么。

\begin{center}\rule{0.5\linewidth}{0.5pt}\end{center}

\subsection{十四、练习题答案}\label{ux5341ux56dbux7ec3ux4e60ux9898ux7b54ux6848}

\textbf{1.} NH\textsubscript{3} + H\textsubscript{2}O \衡 NH\textsubscript{4}\textsuperscript{+} + OH\textsuperscript{-}。共轭酸碱对：(NH\textsubscript{4}\textsuperscript{+}/NH\textsubscript{3}) 和 (H\textsubscript{2}O/OH\textsuperscript{-})。

\textbf{2.} Kb(Ac\textsuperscript{-}) = Kw/Ka(HAc) = 1.0×10\textsuperscript{-14}/1.75×10\textsuperscript{-5} = 5.71×10\textsuperscript{-10}。

\textbf{3.}
(a) 错误。强酸的共轭碱极弱(如 Cl\textsuperscript{-} 是极弱碱)。
(b) 错误。缓冲溶液的 pH 不是恒定不变，而是加入少量酸碱后变化很小。
(c) 错误。37°C 时中性 pH \约等于 6.81(因为 Kw = 2.42×10\textsuperscript{-14})。
(d) 正确。{[}H\textsuperscript{+}{]} = √(Ka\textsubscript{1}·Ka\textsubscript{2}) 与浓度无关。

\textbf{4.} {[}OH\textsuperscript{-}{]} = √(Kb·c) = √(1.77×10\textsuperscript{-5} × 0.10) = 1.33×10\textsuperscript{-3} M。pOH = 2.88，pH = 11.12。

\textbf{5.} CO\textsubscript{3}\textsuperscript{2-} 是二元碱：Kb\textsubscript{1} = Kw/Ka\textsubscript{2} = 2.14×10\textsuperscript{-4}，Kb\textsubscript{2} = Kw/Ka\textsubscript{1} = 2.24×10\textsuperscript{-8}。Kb\textsubscript{1} ≫ Kb\textsubscript{2}，按一元碱处理：{[}OH\textsuperscript{-}{]} = √(Kb\textsubscript{1}·c) = √(2.14×10\textsuperscript{-4} × 0.20) = 6.5×10\textsuperscript{-3} M。pOH = 2.19，pH = 11.81。

\textbf{6.} lg({[}Ac\textsuperscript{-}{]}/{[}HAc{]}) = pH − pKa = 5.00 − 4.76 = 0.24。{[}Ac\textsuperscript{-}{]}/{[}HAc{]} = 10\^{}0.24 = 1.74。

\textbf{7.} 拉平效应是指溶剂将不同强度的酸''拉平''到同一水平的现象。在水中，HClO\textsubscript{4}、HCl、HNO\textsubscript{3} 都把质子完全转移给水生成 H\textsubscript{3}O\textsuperscript{+}，因此在水中它们都是''强酸''，无法区分强弱。

\textbf{8.} Na\textsuperscript{+} 是硬酸，F\textsuperscript{-} 是硬碱，I\textsuperscript{-} 是软碱。硬酸倾向与硬碱结合，所以 NaF 更稳定。

\textbf{9.} c/Ka = 139 \textless{} 400，此时 α = 7.8\% \textgreater{} 5\%，简化公式的误差约 9\%。应使用精确法(二次方程)：x\textsuperscript{2} + Ka·x − Ka·c = 0。

\textbf{10.} pH = pKa\textsubscript{1} + lg({[}HCO\textsubscript{3}\textsuperscript{-}{]}/{[}H\textsubscript{2}CO\textsubscript{3}{]}) = 6.10 + lg(24/1.2) = 6.10 + 1.30 = 7.40 ✓

\begin{center}\rule{0.5\linewidth}{0.5pt}\end{center}

\textbf{11.} 混合弱酸中，Ka(HF) ≫ Ka(HAc)，HAc 的电离被抑制。近似认为 {[}H\textsuperscript{+}{]} \约等于 √(Ka(HF)·c) = √(7.2×10\textsuperscript{-4} × 0.10) = 8.5×10\textsuperscript{-3}，pH \约等于 2.07。

\textbf{12.} 共离子效应：加入 NaAc 后，{[}Ac\textsuperscript{-}{]} 增大，平衡左移，{[}H\textsuperscript{+}{]} 减小。
{[}H\textsuperscript{+}{]} = Ka·{[}HAc{]}/{[}Ac\textsuperscript{-}{]} = 1.75×10\textsuperscript{-5} × 0.10/1.0 = 1.75×10\textsuperscript{-6} M。pH = 5.76。
α = {[}H\textsuperscript{+}{]}/c = 1.75×10\textsuperscript{-6}/0.10 = 0.00175\%(从原来的 1.3\% 降至 0.002\%)。

\textbf{13.} {[}HAc{]} = {[}Ac\textsuperscript{-}{]} = c/2。加入 0.01 mol HCl 后：{[}Ac\textsuperscript{-}{]} = c/2 − 0.01，{[}HAc{]} = c/2 + 0.01。
pH = pKa + lg((c/2−0.01)/(c/2+0.01))。若 c = 1.0 M：pH = 4.76 + lg(0.49/0.51) = 4.76 − 0.017 = 4.74。ΔpH = 0.02。

\textbf{14.} H\textsubscript{3}PO\textsubscript{4} 三步电离，Ka\textsubscript{1} ≫ Ka\textsubscript{2} ≫ Ka\textsubscript{3}：
{[}H\textsuperscript{+}{]} \约等于 [H~2~PO~4~^-^] \约等于 √(Ka\textsubscript{1}·c) = √(7.1×10\textsuperscript{-3} × 0.10) = 2.66×10\textsuperscript{-2} M
{[}HPO\textsubscript{4}\textsuperscript{2-}{]} \约等于 Ka\textsubscript{2} = 6.3×10\textsuperscript{-8} M(与 c 无关)
{[}PO\textsubscript{4}\textsuperscript{3-}{]} \约等于 Ka\textsubscript{2}·Ka\textsubscript{3}/{[}H\textsuperscript{+}{]} = 6.3×10\textsuperscript{-8} × 4.2×10\textsuperscript{-13}/2.66×10\textsuperscript{-2} = 1.0×10\textsuperscript{-18} M

\textbf{15.} {[}H\textsuperscript{+}{]} = 10\textsuperscript{-3}·\textsuperscript{00} = 1.0×10\textsuperscript{-3} M。Ka = {[}H\textsuperscript{+}{]}\textsuperscript{2}/(c−{[}H\textsuperscript{+}{]}) = (1.0×10\textsuperscript{-3})\textsuperscript{2}/(0.10−0.001) = 1.01×10\textsuperscript{-5}。α = 1.0×10\textsuperscript{-3}/0.10 = 1.0\%。

\textbf{16.} pH = 9.25 = pKa(NH\textsubscript{4}\textsuperscript{+})，所以 {[}NH\textsubscript{3}{]}/{[}NH\textsubscript{4}\textsuperscript{+}{]} = 1:1。
设 c\_total = 0.50 M(任意选择)：{[}NH\textsubscript{3}{]} = {[}NH\textsubscript{4}\textsuperscript{+}{]} = 0.25 M。
NH\textsubscript{4}Cl 质量 = 0.25 × 0.50 × 53.49 = 6.69 g。
浓氨水体积 = 0.25 × 0.50/15 = 8.33 mL。

\textbf{17.} Kb\textsubscript{1} = 8.3×10\textsuperscript{-2}，c = 0.10 M，c/Kb\textsubscript{1} = 1.2 ≪ 400，必须精确计算。
设 {[}OH\textsuperscript{-}{]} = x：x\textsuperscript{2}/(0.10−x) = 0.083 \箭 x\textsuperscript{2} + 0.083x − 0.0083 = 0
x = (−0.083 + √(0.083\textsuperscript{2} + 4×0.0083))/2 = (−0.083 + 0.201)/2 = 0.059 M
pOH = 1.23，pH = 12.77。

\textbf{18.} 在液氨中，NH\textsubscript{3} 是溶剂。HCl 和 HAc 都能把质子转移给 NH\textsubscript{3} 生成 NH\textsubscript{4}\textsuperscript{+}，但 NH\textsubscript{3} 的碱性不如 H\textsubscript{2}O 强(作为质子受体)，所以两种酸都只部分电离，表现为弱酸。

\textbf{19.} NaNH\textsubscript{2} + NH\textsubscript{3} \箭 NH\textsubscript{2}\textsuperscript{-} + NH\textsubscript{4}\textsuperscript{+}(NaNH\textsubscript{2} 是强碱，完全夺取 NH\textsubscript{3} 的质子)
NH\textsubscript{4}Cl + NH\textsubscript{3} \箭 NH\textsubscript{3} + NH\textsubscript{4}\textsuperscript{+} + Cl\textsuperscript{-}(NH\textsubscript{4}Cl 是强酸，完全给出质子给 NH\textsubscript{3})

\begin{center}\rule{0.5\linewidth}{0.5pt}\end{center}

\textbf{20.} pH = 7.21 时，主要缓冲对是 H\textsubscript{2}PO\textsubscript{4}\textsuperscript{-}/HPO\textsubscript{4}\textsuperscript{2-}(pKa\textsubscript{2} = 7.20)。
lg({[}HPO\textsubscript{4}\textsuperscript{2-}{]}/{[}H\textsubscript{2}PO\textsubscript{4}\textsuperscript{-}{]}) = 7.21 − 7.20 = 0.01 \箭 比例 = 1.02:1。

H\textsubscript{3}PO\textsubscript{4} + Na\textsubscript{3}PO\textsubscript{4} \箭 3 NaH\textsubscript{2}PO\textsubscript{4}(完全反应后)
但我们需要 H\textsubscript{2}PO\textsubscript{4}\textsuperscript{-}/HPO\textsubscript{4}\textsuperscript{2-} 缓冲，所以不能直接混合。

更实际的方案：H\textsubscript{3}PO\textsubscript{4} + 2Na\textsubscript{3}PO\textsubscript{4} \箭 3Na\textsubscript{2}HPO\textsubscript{4}(控制比例)。
或 H\textsubscript{3}PO\textsubscript{4} + NaOH \箭 NaH\textsubscript{2}PO\textsubscript{4} / Na\textsubscript{2}HPO\textsubscript{4} 缓冲。

\textbf{21.} 三个 Ka 值差 10\textsuperscript{5}\textasciitilde10\textsuperscript{5} 倍，分布系数曲线在 pKa\textsubscript{1}、pKa\textsubscript{2}、pKa\textsubscript{3} 处交叉。
- pH \textless{} 4.7：H\textsubscript{3}PO\textsubscript{4} 主导
- 4.7 \textless{} pH \textless{} 9.8：H\textsubscript{2}PO\textsubscript{4}\textsuperscript{-} 主导
- 9.8 \textless{} pH \textless{} 12.4：HPO\textsubscript{4}\textsuperscript{2-} 主导
- pH \textgreater{} 12.4：PO\textsubscript{4}\textsuperscript{3-} 主导

\textbf{22.} {[}H\textsubscript{2}CO\textsubscript{3}{]} = {[}HCO\textsubscript{3}\textsuperscript{-}{]}/10\^{}(pH−pKa\textsubscript{1}) = 24/10\^{}(7.40−6.10) = 24/20 = 1.2 mmol/L。
{[}CO\textsubscript{2}{]} = 1.2 mmol/L。
PCO\textsubscript{2} = {[}CO\textsubscript{2}{]}/0.030 = 1.2/0.030 = 40 mmHg(正常值 35-45 mmHg)。

\textbf{23.} 半中和点(12.50 mL)时 pH = pKa = 4.00，所以 Ka = 10\textsuperscript{-4}·\textsuperscript{00} = 1.0×10\textsuperscript{-4}。
等当点(25.00 mL)时生成 NaA，pH = 8.00 \箭 pOH = 6.00 \箭 [OH^-^] = 10\textsuperscript{-6}。
Kb = {[}OH\textsuperscript{-}{]}\textsuperscript{2}/c = (10\textsuperscript{-6})\textsuperscript{2}/0.050 = 2.0×10\textsuperscript{-11}。
Ka = Kw/Kb = 10\textsuperscript{-14}/2.0×10\textsuperscript{-11} = 5.0×10\textsuperscript{-4}。
(两组数据给出不同 Ka，说明实验数据有误差，实际取平均或更精确处理。)

\textbf{24.}
(a) AgCl 溶于浓氨水：AgCl + 2NH\textsubscript{3} \箭 [Ag(NH~3~)~2~]\textsuperscript{+} + Cl\textsuperscript{-}，K = Ksp·Kf = 1.8×10\textsuperscript{-10} × 1.1×10\textsuperscript{7} = 2.0×10\textsuperscript{-3}(较大会溶解)。AgI 不溶：K = Ksp·Kf = 8.5×10\textsuperscript{-17} × 1.1×10\textsuperscript{7} = 9.4×10\textsuperscript{-10}(极小不溶)。

\begin{enumerate}
\def\labelenumi{(\alph{enumi})}
\setcounter{enumi}{1}
\item
  HgS 极难溶(Ksp \约等于 10\textsuperscript{-52})，浓 HNO\textsubscript{3} 和王水无法提供足够的驱动力。但 Na\textsubscript{2}S 溶液中 S\textsuperscript{2-} 浓度高，与 Hg\textsuperscript{2+} 形成更稳定的 {[}HgS\textsubscript{2}{]}\textsuperscript{2-} 配离子，使溶解平衡右移。
\item
  Hg\textsuperscript{2+} 是软酸，F\textsuperscript{-} 是硬碱(不稳定)，Cl\textsuperscript{-} 是交界碱(较稳定)。HgF\textsubscript{2} 是离子化合物(强电解质)，HgCl\textsubscript{2} 是共价化合物(弱电解质)。
\end{enumerate}

\textbf{25.}
(a) 未加 NaOH：0.10 M HAc，{[}H\textsuperscript{+}{]} = √(Ka·c) = √(1.75×10\textsuperscript{-5} × 0.10) = 1.32×10\textsuperscript{-3} M，pH = 2.88。

\begin{enumerate}
\def\labelenumi{(\alph{enumi})}
\setcounter{enumi}{1}
\item
  加入 50 mL NaOH：半中和点，{[}Ac\textsuperscript{-}{]} = {[}HAc{]}，pH = pKa = 4.76。
\item
  加入 100 mL NaOH：等当点，生成 0.050 M NaAc。{[}OH\textsuperscript{-}{]} = √(Kb·c) = √(5.71×10\textsuperscript{-10} × 0.050) = 1.69×10\textsuperscript{-5} M。pOH = 4.77，pH = 9.23。
\item
  加入 150 mL NaOH：NaOH 过量 50 mL。{[}OH\textsuperscript{-}{]} = 0.10 × 50/250 = 0.020 M。pOH = 1.70，pH = 12.30。
\item
  等当点 pH = 9.23，应选酚酞(变色范围 8.2\textsubscript{10.0)作为指示剂。甲基橙(3.1}4.4)和甲基红(4.4\textasciitilde6.2)的变色范围都在酸性区域，不适合。
\end{enumerate}
